  % ── Ⅲ. 결 론 ────────────────────────────────────────────
  \chapter{결\quad 론}
  \section{논의 (Discussion)}
  결론은 일반적으로 1개의 장으로 구성하며, 연구결과에 대한 논의, 결론 및 제언 등의 내용을 기술한다.

  \subsection{수식 유효성 조판 검증}
  수식은 지침대로 본문 중앙 정렬 인쇄되며, 수식 번호는 우측 정렬 괄호 인덱싱 처리됩니다.
  \begin{equation}
    Y = ax + b
  \end{equation}

  연구질문에 대한 답변의 경우 현재형 시제를 취하고, 개별 귀결 데이터 수치는 과거형 시제를 취하여 엄격히 기술을 분리합니다.